
\documentclass{article}
\usepackage{colortbl}
\usepackage{makecell}
\usepackage{multirow}
\usepackage{supertabular}

\begin{document}

\newcounter{utterance}

\centering \large Interaction Transcript for game `referencegame', experiment `letter\_grids', episode 10 with reflexion{-}Qwen3.5{-}9B.
\vspace{24pt}

{ \footnotesize  \setcounter{utterance}{1}
\setlength{\tabcolsep}{0pt}
\begin{supertabular}{c@{$\;$}|p{.15\linewidth}@{}p{.15\linewidth}p{.15\linewidth}p{.15\linewidth}p{.15\linewidth}p{.15\linewidth}}
    \# & $\;$A & \multicolumn{4}{c}{Game Master} & $\;\:$B\\
    \hline

    \theutterance \stepcounter{utterance}  
    & & \multicolumn{4}{p{0.6\linewidth}}{
        \cellcolor[rgb]{0.9,0.9,0.9}{
            \makecell[{{p{\linewidth}}}]{
                \texttt{\tiny{[P1$\langle$GM]}}
                \texttt{You have attempted this same task before and did not succeed. These are the plans you wrote after those attempts.} \\
\texttt{Trial \#1: New plan: My previous description failed because it incorrectly identified the "bottom{-}left column" as completely filled (it is only filled in the first three rows) and mischaracterized the empty area as a "2x4 area" (the empty region is actually a 3x4 block spanning rows 2{-}4 and columns 2{-}5). To succeed, I must strictly adhere to the visual evidence: describe the top row as fully filled, the leftmost column as filled only for the first three rows, and the third row as fully filled. I will explicitly state that the remaining cells (rows 2{-}4, columns 2{-}5) are empty to ensure the description uniquely distinguishes the target from the distractors, which have different patterns of filled rows and columns.} \\
\\ 
\texttt{Trial \#2: New plan: My previous description failed because it incorrectly stated that the leftmost column is filled only in the first three rows; in reality, the leftmost column is completely filled (rows 1{-}5). Additionally, I misidentified the empty region as spanning rows 2{-}4, when it actually spans rows 2{-}5. To succeed, I must strictly adhere to the visual evidence: describe the top row (row 1) and the third row (row 3) as completely filled, the leftmost column as completely filled, and explicitly state that the remaining cells (rows 2, 4, and 5, columns 2 through 5) are empty. This precise distinction will differentiate the target from Distractor 1 (which has a filled second row) and Distractor 2 (which has a filled fifth row).} \\
\\ 
\texttt{Trial \#3: New plan: My previous attempts failed because I generated verbose explanations and internal monologues instead of providing only the required "Expression: " tag followed immediately by the description. I also made factual errors regarding the grid contents, specifically misidentifying the fill status of Row 3 and the extent of the empty region. To succeed, I will strictly output only the tag "Expression: " followed by a concise, accurate description: "The target grid has a completely filled top row (row 1) and a completely filled third row (row 3), with the leftmost column (column 1) fully filled from top to bottom. All other cells are empty." This description correctly distinguishes the target from Distractor 1 (which has a filled row 2) and Distractor 2 (which has a filled row 5).} \\
\\ 
\\ 
\texttt{You are given three grids, where each of them is 5 by 5 in size.} \\
\texttt{Grids have empty cells marked with "▢" and filled cells marked with "X".} \\
\texttt{Your task is to generate a referring expression that best describes the target grid while distinguishing it from the two other distractor grids.} \\
\texttt{The first grid is the target grid, and the following two grids are the distractors.} \\
\\ 
\texttt{Target grid:} \\
\\ 
\texttt{X X X X X} \\
\texttt{□ □ □ □ □} \\
\texttt{X X X X □} \\
\texttt{X □ □ □ □} \\
\texttt{X □ □ □ □} \\
\\ 
\texttt{Distractor grid 1:} \\
\\ 
\texttt{X X X X X} \\
\texttt{X □ □ □ □} \\
\texttt{X X X X X} \\
\texttt{X □ □ □ □} \\
\texttt{X X X X X} \\
\\ 
\texttt{Distractor grid 2:} \\
\\ 
\texttt{X X X X X} \\
\texttt{X □ □ □ □} \\
\texttt{X □ □ □ □} \\
\texttt{X □ □ □ □} \\
\texttt{X X X X X} \\
\\ 
\texttt{Instruction: Describe the target grid.} \\
\texttt{Generate the referring expression starting with the tag "Expression: " for the given target grid. Omit any other text.} \\
            }
        }
    }
    & \\ \\

    \theutterance \stepcounter{utterance}  
    & \multicolumn{4}{p{0.6\linewidth}}{
        \cellcolor[rgb]{0.9,0.9,0.9}{
            \makecell[{{p{\linewidth}}}]{
                \texttt{\tiny{[P1$\rangle$GM]}}
                \texttt{Expression: The target grid has a completely filled top row (row 1) and a completely filled third row (row 3), with the leftmost column (column 1) fully filled from top to bottom. All other cells are empty.} \\
            }
        }
    }
    & & \\ \\

    \theutterance \stepcounter{utterance}  
    & & & \multicolumn{2}{p{0.3\linewidth}}{
        \cellcolor[rgb]{0.9,0.9,0.9}{
            \makecell[{{p{\linewidth}}}]{
                \texttt{\tiny{[GM$|$GM]}}
                \texttt{Expression: The target grid has a completely filled top row (row 1) and a completely filled third row (row 3), with the leftmost column (column 1) fully filled from top to bottom. All other cells are empty.} \\
            }
        }
    }
    & & \\ \\

    \theutterance \stepcounter{utterance}  
    & & \multicolumn{4}{p{0.6\linewidth}}{
        \cellcolor[rgb]{0.9,0.9,0.9}{
            \makecell[{{p{\linewidth}}}]{
                \texttt{\tiny{[GM$\rangle$P2]}}
                \texttt{You have attempted this same task before and did not succeed. These are the plans you wrote after those attempts.} \\
\texttt{Trial \#1: New plan: I incorrectly identified the "First" grid because I focused on the filled top and bottom{-}left columns but failed to verify the third row. The expression requires the *third* row to be completely filled, which the "First" grid lacks (it has empty cells in the third row). The "Third" grid matches all criteria: a full top row, a full bottom{-}left column, a full third row, and an empty bottom{-}right 2x4 area. Next time, I will systematically check each grid against all three specific conditions (top row, third row, and bottom{-}left column) before making a selection, ensuring no condition is overlooked.} \\
\\ 
\texttt{Trial \#2: New plan: I incorrectly selected the "First" grid because I focused on the filled top row and the leftmost column pattern but failed to verify the third row condition. The expression explicitly requires the third row to be completely filled, which the "First" grid lacks (it has empty cells in the third row). The "Third" grid matches all criteria: a full top row, a full third row, and the specific leftmost column pattern. Next time, I will systematically check each grid against all three specific conditions (top row, third row, and bottom{-}left column) before making a selection, ensuring no condition is overlooked.} \\
\\ 
\\ 
\texttt{You are given three grids, where each of them is 5 by 5 in size.} \\
\texttt{Grids have empty cells marked with "▢" and filled cells marked with "X".} \\
\texttt{You are also given a referring expression that describes one of the given grids.} \\
\texttt{Your task is to select the grid that matches the given referring expression.} \\
\texttt{Generate only the number (in text) of the grid that matches the given expression by selecting first, second, or third.} \\
\\ 
\texttt{First:} \\
\\ 
\texttt{X X X X X} \\
\texttt{X □ □ □ □} \\
\texttt{X X X X X} \\
\texttt{X □ □ □ □} \\
\texttt{X X X X X} \\
\\ 
\texttt{Second:} \\
\\ 
\texttt{X X X X X} \\
\texttt{□ □ □ □ □} \\
\texttt{X X X X □} \\
\texttt{X □ □ □ □} \\
\texttt{X □ □ □ □} \\
\\ 
\texttt{Third:} \\
\\ 
\texttt{X X X X X} \\
\texttt{X □ □ □ □} \\
\texttt{X □ □ □ □} \\
\texttt{X □ □ □ □} \\
\texttt{X X X X X} \\
\\ 
\texttt{Expression: The target grid has a completely filled top row (row 1) and a completely filled third row (row 3), with the leftmost column (column 1) fully filled from top to bottom. All other cells are empty.} \\
\texttt{Question: Which grid does the expression refer to?} \\
\texttt{Start with the tag "Answer: ", followed by your selection. Omit any other text.} \\
            }
        }
    }
    & \\ \\

    \theutterance \stepcounter{utterance}  
    & & & \multicolumn{4}{p{0.6\linewidth}}{
        \cellcolor[rgb]{0.9,0.9,0.9}{
            \makecell[{{p{\linewidth}}}]{
                \texttt{\tiny{[GM$\langle$P2]}}
                \texttt{Answer: First} \\
            }
        }
    }
     \\ \\

    \theutterance \stepcounter{utterance}  
    & & & \multicolumn{2}{p{0.3\linewidth}}{
        \cellcolor[rgb]{0.9,0.9,0.9}{
            \makecell[{{p{\linewidth}}}]{
                \texttt{\tiny{[GM$|$GM]}}
                \texttt{first} \\
            }
        }
    }
    & & \\ \\

    \theutterance \stepcounter{utterance}  
    & \multicolumn{4}{p{0.6\linewidth}}{
        \cellcolor[rgb]{0.95,0.95,0.95}{
            \makecell[{{p{\linewidth}}}]{
                \texttt{\tiny{[P2$\rightarrow$P1]}}
                \texttt{Post{-}episode Reflexion (not sent to GM):} \\
\texttt{New plan: My previous attempts failed because I incorrectly identified the extent of the filled leftmost column and the empty region. Upon re{-}examining the target grid, I see that the leftmost column (column 1) is completely filled from row 1 to row 5, not just the first three rows. Additionally, the empty region spans rows 2, 4, and 5 in columns 2 through 5, not a continuous block from row 2 to 4. To succeed, I will generate a concise expression that explicitly states: "The target grid has a completely filled top row (row 1), a completely filled third row (row 3), and a completely filled leftmost column (column 1). All other cells (rows 2, 4, and 5 in columns 2{-}5) are empty." This precise description will correctly distinguish the target from Distractor 1 (which has a filled row 2) and Distractor 2 (which has a filled row 5).} \\
            }
        }
    }
    & & \\ \\

    \theutterance \stepcounter{utterance}  
    & & & \multicolumn{4}{p{0.6\linewidth}}{
        \cellcolor[rgb]{0.95,0.95,0.95}{
            \makecell[{{p{\linewidth}}}]{
                \texttt{\tiny{[P1$\rightarrow$P2]}}
                \texttt{Post{-}episode Reflexion (not sent to GM):} \\
\texttt{New plan: I failed because I stopped checking after verifying the top row and the leftmost column, assuming the first grid was correct without confirming the third row condition. The "First" grid has empty cells in the third row, violating the requirement for a completely filled third row. The "Third" grid is the only one that satisfies all three conditions: a full top row, a full third row, and a fully filled leftmost column. Next time, I will explicitly verify the third row of the "First" grid before selecting it, and if it fails, immediately select the "Third" grid as the correct match.} \\
            }
        }
    }
     \\ \\

\end{supertabular}
}

\end{document}
